\clearpage
\section{辅助工具的使用}

本次实验用豆包与 Kimi 参与设计：需求到模块的切分、六类 UML 的样例代码、界面版式与控件方案。三条使用原则：先有自己的判断再问模型（把它当能立刻反驳的同伴，而不是答案来源）；只信可验证的中间产物（能渲染的代码、能逐条核对的表）；输出必须改写进自己的结构，不做“看着不错就贴进去”。

\subsection{使用记录}

五次影响设计的对话见表~\ref{tab:ai}，每次记录场景、生成结果与问题、修改后的结果。

\begin{center}
\captionof{table}{辅助工具使用记录（场景 / 生成结果与问题 / 修改后）}\label{tab:ai}
\zihao{-5}
\renewcommand{\arraystretch}{1.15}
\begin{tabularx}{\textwidth}{@{}p{2.95cm}Yp{5.15cm}@{}}
\toprule
使用场景 & 生成结果与问题 & 修改优化后的结果 \\
\midrule
\textbf{需求到模块}：把 19 条需求与异常清单给它查漏 & 指出“分类管理无模块承接”“定时调度无归属”两处缺口；但 6 模块方案切得过粗，审核要读写商品却独立成模块，也没考虑日志、通知这类横切能力 & 调成 8 模块：分类并入商品，新增日志统计与公共支撑。模型能给覆盖面，给不了“边界切在哪”（表~\ref{tab:module}） \\
\textbf{类图代码}：先列实体清单，让它补属性、方法与多重性 & 24 个类 8 个包，还凭常识加了 Order 实体，与“不保存独立交易事实”冲突；宽高比 4.5:1，放到 A4 上已看不清 & 删 Order；服务类移到组件图；只留 14 个实体与枚举，补组合关系与幂等注释，宽高比降到 1.9:1（\figref{fig:uml1}） \\
\textbf{“修改待审”建模}：把 REQ-017 原文与状态表给它 & 先建议新增“修改待审”状态；追问“驳回后怎么回去”才改为组合状态。代码漏了内部初始转移，且默认“审核期间不能预约”，与需求相反 & 保留组合状态并补三条内部转移；待审内容落成 ItemRevision，驳回即丢弃（\figref{fig:uml5}） \\
\textbf{时序图异常路径}：我写初版，让它检查漏了哪些分支 & 补出“失败也写日志”“通知应在事务提交后发”；但给出“先加锁再查 requestId”，重复请求会排队等锁 & 改为“先查 requestId、未命中再开事务加锁”，三种失败分支全部画进 alt（\figref{fig:uml6}） \\
\textbf{界面布局与控件}：说明“我的商品”要表达状态与动作、后台要排列表与详情 & 三种布局都推荐表格；后台建议左右分栏、驳回原因用下拉框；但 390 宽用表格要横向滚动，下拉框不利于统计，灰阶文字对比度也不足 & 改卡片列表 + 状态标签 + 按状态绑定的按钮；驳回原因改标签组；灰阶文字统一压深（\figref{fig:ui34}） \\
\bottomrule
\end{tabularx}
\end{center}

\subsection{把 AI 放进工作流}

上表的共同点是：模型擅长的部分都能被机器立刻判真假（puml 能否渲染、表格能否逐条核对、截图能否看清），出错的部分都来自它没有的信息（业务约束、版式习惯）。据此可以把流程重排成“模型生成、机器校验、人只做检查校对”，本实验实际跑的五步见表~\ref{tab:aiflow}。

\begin{center}
\begin{minipage}{\textwidth}\centering
\captionof{table}{“模型生成 + 机器校验 + 人工校对”工作流}\label{tab:aiflow}
\zihao{-5}
\renewcommand{\arraystretch}{1.12}
\begin{tabularx}{\textwidth}{@{}p{1.55cm}Yp{4.3cm}p{2.6cm}@{}}
\toprule
环节 & 模型负责 & 人负责 & 可验证的产物 \\
\midrule
① 定骨架 & 不参与 & 从需求划模块、列图清单、写下每张图要表达的结论 & 模块表、图清单 \\
② 出产物 & 生成 puml / 原型 HTML / 表格骨架 & 给约束与字段，点出不合业务的地方 & 能渲染的代码 \\
③ 机器校验 & 不参与 & 渲染、编译、逐条比对需求编号 & 出图、出 PDF \\
④ 检查校对 & 可按清单自查（仅作提醒） & 三张清单：需求覆盖、图形语义、界面一致 & 修订 diff \\
⑤ 入库 & 不参与 & 分次提交并注明 AI 生成或人工修改 & git log \\
\bottomrule
\end{tabularx}
\end{minipage}
\end{center}

“人只做检查校对”成立需要三个前提。其一，每个产物都有机器可判定的通过标准，本实验中是“puml 能渲染且中文不乱码”“PDF 无溢出且图表编号连续”“19 条需求每条都能指到图、表或界面”；没有通过标准的产物（例如“这个架构好不好”）无法交给模型。其二，需求可编号，覆盖关系才能逐条核对，这也是实验一把需求编到 REQ-019 的直接收益。其三，有回退手段，模型整段重写是常事，分次提交后只需撤掉某一次（第~\ref{sec:git} 节）。

剩下唯一不能外包的是检查本身。本报告的文字也用了模型扩写，做法是先自己写下一节的一句话结论与三条要点，再让模型扩成段落，最后拿结论去核对图、截图与需求编号；模型自行添加的判断一律回到产物里验证，验证不了就删。这一轮删掉三处读起来合理但没有依据的句子，例如初稿曾写“审核期间用户不能预约”，与 REQ-017 恰好相反。

\subsection{优缺点与小结}

优点是查漏快、语法模板可靠、概念映射有新角度（把业务描述映射成 UML 组合状态）、方案对比高效。不足是会凭常识造需求、不懂业务约束（认为审核期间不可预约）、没有版式意识、存在顺序与性能盲区（把加锁放在幂等查重之前）、视觉细节不可靠。

由此得到三点判断。其一，这类工具适合做“陪练”而不是“代笔”：设计的核心判断依赖对需求的熟悉程度，模型没有这个上下文，但当先有判断时它能立刻给出反驳与替代方案。其二，它擅长形式、不擅长约束：PlantUML 语法与布局模板表现很好，而“不做支付所以不该有订单”“审核期间仍可预约”这类约束只能来自需求文档。其三，要让输出可被证伪：能立刻渲染的图、能逐条核对的表价值最高。为了让模型理解“修改待审”，我被迫把 REQ-017 逐字读了三遍，才发现自己原以为“审核期间不能预约”——写清问题的过程本身就是一次需求复核。
